\documentclass[a4paper, 11pt]{ltjsarticle}

% --- 日本語環境とフォント設定 ---
\usepackage{luatexja-fontspec}
\usepackage[ipaex]{luatexja-preset}

% --- 必要なパッケージ ---
\usepackage{graphicx}
\usepackage{listings}
\usepackage{amsmath}
\usepackage{geometry}
\usepackage{hyperref}
\usepackage{booktabs}      % 綺麗な表のため
\usepackage{multirow}    % 表のセル結合

% --- ページレイアウト ---
\geometry{
  left=25mm,
  right=25mm,
  top=30mm,
  bottom=30mm
}

% --- タイトル情報 ---
\title{短期株価予測におけるファンダメンタルズ指標の有効性検証\\（V2公正比較実験）}
\author{（あなたの氏名）}
\date{\today}

% ======================================
% --- ドキュメント開始 ---
% ======================================
\begin{document}

\maketitle

% ======================================
\section{試験概要}
% ======================================
初期分析（V1）において、純粋な財務指標のみを用いた短期（3ヶ月後）株価予測モデルの
正解率が50\%（ランダム）程度にとどまることが確認された。

本試験の目的は、このベースラインモデルに対し、株価と財務指標を組み合わせた
「バリュー指標」を追加することで、予測精度が改善するかを検証することである。
分析対象は、赤字企業なども含む「ごちゃまぜデータ」全体
（\texttt{dropna(subset=['target'])}ルールで作成）とした。

以下の2つのアプローチを、同一のデータ作成ルールと分析手法を用いて公正に比較した。
\begin{itemize}
    \item \textbf{V2 (PER/PBR版):} 人間が定義した「完成品」の指標（\texttt{PER}, \texttt{PBR}）を追加する。
    \item \textbf{V2\_Close (部品版):} 「完成品」の計算に使われる「部品」
    （\texttt{Close}, \texttt{発行済株式総数}）を直接追加する。
\end{itemize}

% ======================================
\section{実験1：V2 (PER/PBR版) の結果と考察}
% ======================================
\subsection*{1. 結果}
\begin{itemize}
    \item \textbf{正解率 (Accuracy): 0.5094 (50.9\%)}
    \item \textbf{トップ3重要特徴量:}
    \begin{enumerate}
        \item \texttt{Profit} (純利益) : 0.339
        \item \texttt{ROE} (自己資本利益率) : 0.273
        \item \texttt{OperatingProfit} (営業利益) : 0.134
    \end{enumerate}
    \item \textbf{追加特徴量の重要度:} \texttt{PER} (0.018), \texttt{PBR} (0.019) ともに最下位クラス。
\end{itemize}

\subsection*{2. 考察}
正解率は50.9\%と、V1のベースライン（約50\%）とほぼ変わらず、ランダムな予測と
同等であった。

モデルは、追加された\texttt{PER}や\texttt{PBR}を予測にほとんど使用せず
（重要度最下位）、V1と同様に\texttt{Profit}や\texttt{ROE}を重視しようとした。
しかし、赤字企業も含む「ごちゃまぜデータ」においては、それらの財務指標も
有効なシグナルとはならなかった。

この結果から、「ごちゃまzeデータ」の短期予測において、
伝統的なバリュー指標（PER/PBR）を追加するアプローチは有効ではないことが示された。

% ======================================
\section{実験2：V2\_Close (部品版) の仮説・結果・考察}
% ======================================
\subsection*{1. 仮説}
実験1において、「完成品」の\texttt{PER}/\texttt{PBR}指標は予測に効かなかった。
そこで、「完成品」を与えるよりも、その「部品」である\texttt{Close}（日々の株価）や
\texttt{発行済株式総数}を直接モデルに与えた方が、モデル自身がより有効なパターン
（例えば「時価総額」のような概念）を発見し、予測精度が改善するのではないか、
という仮説を立てた。

\subsection*{2. 結果}
\begin{itemize}
    \item \textbf{正解率 (Accuracy): 0.5660 (56.6\%)}
    \item \textbf{トップ3重要特徴量:}
    \begin{enumerate}
        \item \texttt{TotalAssets} (総資産) : 0.217
        \item \texttt{Profit} (純利益) : 0.191
        \item \texttt{OperatingProfit} (営業利益) : 0.177
    \end{enumerate}
    \item \textbf{追加特徴量の重要度:} \texttt{Issued Shares} (0.096, 4位), \texttt{Close} (0.046, 8位) となり、
    V2のPER/PBRよりは高い重要度を示した。
\end{itemize}

\subsection*{3. 考察}
仮説はわずかに正しかった。正解率は50.9\%から56.6\%へと微増し、「完成品」よりも
「部品」を与えた方が、モデルがより多くのシグナルを抽出できたことが示唆された。
特徴量の重要度も分散し、モデルが特定の指標に頼るのではなく、
複数の「部品」を組み合わせてパターンを探ろうと試みた様子がうかがえる。

\subsection*{4. 総合結論}
本比較実験の最大の発見は、どちらのアプローチも正解率は50\%～56\%の範囲にとどまり、
**「ごちゃまzeデータ」の短期予測という土俵では、ファンダメンタルズ指標や
バリュー指標だけでは、強力な予測モデルを構築するのは困難である**、
ということである。

これは、田村氏の研究（第3章）の「短期予測ではファンダメンタルズ分析は効きにくい」
という結論と完全に一致する。我々の「ごちゃまzeデータ」の正解率（50\%台）を
引き上げるためには、ファンダメンタルズとは全く異なる種類の武器、すなわち
**テクニカル指標（株価の勢い）**の追加が不可欠であると結論付けられる。

\end{document}